\section{Background and Related Work}
\label{sec:background}

\subsection{Interactive AI Agents and the Network Assumption}
\label{sec:agents}

Interactive AI agents operate as closed-loop controllers that iteratively reason, act, and observe. The ReAct framework \cite{yao2023react} introduced this paradigm: at each step, the agent issues a command, parses the resulting output, and conditions its next action on that observation. Terminal interfaces have emerged as the primary execution channel for such agents. SWE-agent \cite{Yang2024} demonstrated that agents using well-designed command-line interfaces achieve higher task success rates (12.47\% vs. 3.8\%) than agents using generic tool APIs. While Terminal-Bench \cite{merrill2026} formalized evaluation across system administration, scripting, and infrastructure management tasks. A critical assumption pervades the agent literature: the underlying remote shell connection is treated as reliable and low-latency. Because the agent's control loop is tightly coupled to its feedback channel, transport-layer latency and reliability directly govern execution coherence.

\subsection{Remote Shell Protocols}
\label{subsec:protocols}

The resilience of an interactive AI agent depends on the transport protocol underlying its shell session. Each entails different trade-offs among latency, completeness, and connection setup overhead.

\paragraph{SSHv2}
The Secure Shell Protocol version 2 \cite{Ylonen2006RFC4251} is the de facto mechanism for remote terminal access. It operates over TCP, multiplexing encrypted channels over a single connection with ordered, reliable delivery \cite{Bellare2002}. However, SSHv2 inherits TCP's HOL blocking, where a single lost segment stalls all subsequent data until retransmission completes \cite{Scharf2007}. For agents in tight action-observation loops, these stalls manifest as tail latency spikes that may exceed application-layer timeout thresholds. Session establishment requires 2-3 round trips: one for the TCP handshake and one to two for the SSH cryptographic handshake \cite{Ylonen2006RFC4253}, imposing setup latency that compounds with network RTT.

\paragraph{SSH3}
SSH3 \cite{michel2023} redesigns SSH over HTTP/3 and the QUIC transport \cite{Iyengar2021}. By integrating TLS 1.3 into the transport handshake, QUIC reduces session establishment to 1-RTT (or 0-RTT for resumed connections), significantly lowering initial latency compared to SSHv2. While QUIC eliminates cross-stream HOL blocking, a single interactive session still relies on one sequential stream where intra-stream loss requires retransmission. Its practical advantage lies in tighter loss detection and userspace recovery, which may reduce latency variance. As a nascent protocol, its empirical behavior under degraded conditions remains unevaluated.

\paragraph{Mosh}
Mosh \cite{Winstein2012} was designed to address usability failures of SSHv2 over unreliable networks by diverging from the byte-stream paradigm entirely. After an initial SSH session for authentication, Mosh switches to a UDP channel carrying terminal state updates via the State Synchronization Protocol (SSP). SSP treats the terminal as a replicated state machine, transmitting only compact diffs of the latest screen state; lost datagrams are simply superseded by newer ones, eliminating HOL blocking completely. Mosh additionally masks network delay through predictive local echo, rendering keystrokes before server confirmation. This design trades completeness for responsiveness: intermediate output may be elided before delivery. For an agent that relies on parsing complete command output, this elision can withhold output that the agent's next action depends on, such as stack traces or build logs.

\paragraph{Prior Measurement Studies}
Prior measurement studies focus on different dimensions than ours. \cite{Winstein2012} measured Mosh's keystroke response time against SSHv2 under real-world mobile connectivity, reporting substantially lower median keystroke latency over 3G networks compared to SSH. However, their evaluation targeted user-perceived responsiveness under natural roaming, without controlled degradation profiles or tail latency analysis. On the QUIC side, \cite{Langley2017} reported Internet-scale deployment results showing reduced page-load times and rebuffering rates for web and video workloads, while \cite{Kakhki2017} conducted comprehensive cross-environment measurements of QUIC's transport behavior across desktop, mobile, wired, and wireless settings. Both studies evaluated bulk transfer and page-load performance; neither extended their analysis to interactive terminal sessions where per-command latency and its variance, rather than bulk throughput, constitute the governing metrics. On the agent evaluation side, SWE-agent \cite{Yang2024} and Terminal-Bench \cite{merrill2026} benchmarked task success rates and completion times but assumed ideal local or low-latency connections and did not vary network conditions as an independent variable. No prior study has placed SSHv2, SSH3, and Mosh on a common workload suite representative of interactive AI agent workloads under systematically controlled network conditions.


\subsection{The Measurement Gap}
\label{subsec:gap}

These architectural differences produce distinct latency profiles that interact with agent behavior with measurable consequences for agent reliability. SSHv2 guarantees completeness but exposes agents to HOL-induced tail latency under loss. SSH3 accelerates loss recovery and session setup but retains sequential delivery constraints within a single stream. Mosh eliminates transport-layer stalls at the cost of potentially incomplete observations. The extent to which these trade-offs manifest under realistic network impairment, and their quantitative impact on the latency distributions that agent timeout mechanisms must accommodate, has not been empirically established. This study provides that characterization.
